% HSK 电工杯 LaTeX 通用起稿模板 v7.6.0
% 编译：python scripts/render_paper.py final_latex --profile diangong --clean
% 当届规则、页数、封面和指定章节要求始终优先于本通用模板。
\documentclass[12pt,a4paper]{ctexart}

\usepackage[margin=2.54cm]{geometry}
\usepackage{amsmath,amssymb,amsthm,bm}
\usepackage{graphicx}
\usepackage{float}
\usepackage{booktabs}
\usepackage{longtable}
\usepackage{multirow}
\usepackage{listings}
\usepackage{xcolor}
\usepackage{hyperref}
\usepackage{caption}
\usepackage{subcaption}
\usepackage{enumitem}
\usepackage{titlesec}
\usepackage{fancyhdr}
\usepackage{lastpage}
\usepackage[backend=biber,style=numeric,sorting=none]{biblatex}
\addbibresource{references.bib}

\newcommand{\ProblemCode}{A}
\newcommand{\TeamNumber}{XXXX}
\newcommand{\PaperTitle}{电工杯论文标题}

\hypersetup{colorlinks=false,pdfborder={0 0 0}}
\lstset{
    basicstyle=\small\ttfamily,
    breaklines=true,
    frame=single,
    numbers=left,
    numberstyle=\tiny,
    showstringspaces=false
}

\pagestyle{fancy}
\fancyhf{}
\fancyhead[L]{队伍编号 \TeamNumber}
\fancyhead[R]{第 \thepage\ 页 / 共 \pageref{LastPage} 页}
\renewcommand{\headrulewidth}{0.4pt}
\titleformat*{\section}{\Large\bfseries}
\titleformat*{\subsection}{\large\bfseries}
\titleformat*{\subsubsection}{\normalsize\bfseries}

\begin{document}

\begin{titlepage}
\begin{center}
\vspace*{2cm}
{\Huge\bfseries \PaperTitle}\\[2cm]
{\Large 题号：\textbf{\ProblemCode}}\\[1cm]
{\Large 队伍编号：\textbf{\TeamNumber}}\\[3cm]
\end{center}
\end{titlepage}

\section*{摘要}
% 第一段简要交代工程对象、核心矛盾和统一路线；随后按实际小问写模型/机制、决定性数值和直接结论。
% 优先正向说明“建立什么—得到什么—说明什么”，不罗列程序和工作簿。
填写摘要正文。

\noindent\textbf{关键词：}关键词1；关键词2；关键词3；关键词4

\newpage
\tableofcontents
\newpage

\section{问题重述}
\subsection{问题背景}
% 通常一个自然段。由题面工程背景和真实工程矛盾自然引出待解决问题。

\subsection{问题提出}
\noindent\textbf{问题一：}……

\noindent\textbf{问题二：}……
% 按实际小问继续；只转述对象、条件和输出，不提前写公式、算法和结果。

\section{问题分析}
\subsection{问题一的分析}
% 写工程难点、对象关系、关键约束和建模抓手；正式公式和最终结果默认留到模型/结果章节。

\subsection{问题二的分析}
% 按实际小问继续。跨问依赖可以说明，但不重复问题提出。

\section{模型假设}
\begin{enumerate}
    \item ……。
    \item ……。
\end{enumerate}

\section{符号说明}
\begin{table}[H]
    \centering
    \caption{主要符号说明}
    \begin{tabular}{cccc}
        \toprule
        符号 & 含义 & 单位 & 类型 \\
        \midrule
        $x_i$ & 示例变量 & -- & 变量 \\
        $p^{(s)}$ & 第 $s$ 个场景参数 & -- & 参数 \\
        \bottomrule
    \end{tabular}
\end{table}

% preprocessing_decision=not_needed 时删除本章；question_local 放到对应问题；project_level 才保留项目级预处理。
\section{数据预处理}
\subsection{数据问题与处理必要性}
\subsection{处理方法与验证}

\section{问题一模型建立及求解}
\subsection{工程机制与模型建立}
% 从本题工程对象和约束出发，定义变量并推导关键关系。

% 核心模型收束按项目框架中的 required / inline / not_applicable 自适应。
% 下面展示 required 示例；inline 时在上一小节末尾直接给最终关系并删除本小节；not_applicable 时直接删除。
\subsection{核心模型汇总}
\begin{equation}
\left\{
\begin{aligned}
\min\quad & f(x) \\
\text{s.t.}\quad & g_i(x)\le 0,\\
& h_j(x)=0,\\
& x\in\Omega .
\end{aligned}
\right.
\end{equation}

\subsection{模型求解}
% 写实际算法、初始化/参数、终止条件、工程可行性和收敛口径。

\subsection{求解结果}
% 主结果留在本问。核心图表必须在邻近正文显式引用编号，并解释工程量级、约束和机制。
% 例如：图~\ref{fig:q1-main} 给出了……；表~\ref{tab:q1-main} 汇总了……。

% 若存在实质结果深化，使用题目专属标题，例如：
% \subsection{参数扰动下的工程稳定区间}
% \subsection{极端工况压力测试}
% 最后一个结果/深化段自然回答本问，不默认建立“小问结论”。

% 按实际题目继续问题二及后续；每问是否单列“核心模型汇总”由当前模型复杂度决定。

\section{模型的评价与推广}
% 确有工程改进内容时可改为“模型的改进、评价与推广”。
\subsection{模型的优点}
% 数量由实际模型决定；每条必须有结构或验证证据，不检查与缺点数量的大小关系。
\begin{enumerate}
    \item ……。
    \item ……。
    \item ……。
\end{enumerate}

\subsection{模型的缺点}
% 数量由实际限制决定，并说明影响的结果、参数范围或适用边界。
\begin{enumerate}
    \item ……。
    \item ……。
\end{enumerate}

\subsection{模型的推广}
% 选写：说明可迁移工程对象、需替换参数/边界和不可直接迁移条件。

% 默认不设置独立“结论”章；各问直接答案已经在对应“求解结果”中闭环。
% 当届电工杯规则若明确要求总结/结论章节，则按当届规则增加。

\printbibliography[title={参考文献}]

\appendix
\section{主要算法与复现说明}
\section{补充推导与图表}
% 完整代码可独立提交；附录说明输入、功能、输出和复现顺序。

\end{document}